% funding-rate.tex — Bidirectional funding rate with jump adjustment
% Kontrol: prove_cfmm_perpetual_fundingConvergence, prove_cfmm_perpetual_jumpSafety

\section{Funding Rate}\label{sec:funding-rate}

This instrument is perpetual: it has no fixed expiry but settles at liquidity events.
The funding rate maintains mark-index convergence, following the framework of
\cite{angeris2022perpetuals}. It is separate from the borrow rate
(\cref{sec:borrow-rate}) and bidirectional.

\subsection{Mark and Index Prices}

\begin{definition}[Index Price]\label{def:index-price}
The index price is derived from the FCI hook's co-primary state:
\begin{equation}\label{eq:index-price}
  \pindex = \frac{\DeltaPlus}{1 - \DeltaPlus}
\end{equation}
where $\DeltaPlus = \max(0,\, \AT - \ATnull)$ is the realized concentration deviation
from the underlying pool. Computed on the fly from stored $\AT$, $\ThetaSum$, $\PosCount$.
\end{definition}

\begin{definition}[Mark Price]\label{def:mark-price}
The mark price is the CFMM's implied price from virtual reserves:
\begin{equation}\label{eq:mark-price}
  \pmark = \frac{\pl^2}{2}\sqrt{\frac{u}{\Lact}}
\end{equation}
where $u$ is the current virtual concentration reserve and $\Lact$ is
the active liquidity in the current tick range.
\end{definition}

\subsection{Discrete Funding at Hook Callbacks}

Funding is exchanged at each hook callback (\texttt{afterSwap},
\texttt{afterAddLiquidity}, \texttt{afterRemoveLiquidity}), matching
the FCI oracle update frequency. No keeper or per-block updates required.

\begin{definition}[Funding Payment]\label{def:funding-payment}
At callback $n$, the funding rate is:
\begin{equation}\label{eq:funding-rate}
  \rfunding_n = \alpha \cdot \frac{\pmark_n - \pindex_n}{\pindex_n + 1}
\end{equation}
where $\alpha > 0$ is the sensitivity parameter (set at pool initialization).
The denominator $(\pindex + 1)$ normalizes by index level, preventing
extreme rates when $\pindex$ is small.
\end{definition}

\begin{remark}[Bidirectionality]
\begin{itemize}
  \item $\pmark > \pindex$ (protection overpriced): $\rfunding > 0$,
    PLPs pay underwriters. Incentivizes arbitrageurs to sell the perpetual,
    pushing $\pmark$ down.
  \item $\pmark < \pindex$ (protection underpriced): $\rfunding < 0$,
    underwriters pay PLPs. Incentivizes buying the perpetual,
    pushing $\pmark$ up.
\end{itemize}
This bidirectionality is essential for convergence and is why funding must
be in a separate accumulator from the unidirectional borrow rate
(\cref{sec:borrow-rate}).
\end{remark}

\subsection{Jump-Adjusted Replication}

The backtest (ETH/USDC 30bps, 41 days) confirms that $\DeltaPlus$ exhibits
rare, large, discontinuous jumps: $0 \to 0.158$ in a single day.
Between callbacks, a single large swap or JIT event can change the FCI
dramatically. The funding rate includes a jump premium.

\begin{definition}[Jump Premium]\label{def:jump-premium}
When $|\Delta\pindex| > \delta_{\text{jump}}$ at callback $n$
(where $\delta_{\text{jump}}$ is the jump threshold), the funding
rate receives an additive adjustment:
\begin{equation}\label{eq:jump-funding}
  \rfunding_n^{\text{adj}} = \rfunding_n + \lambda \cdot |\Delta\pindex|
\end{equation}
where $\lambda > 0$ is the jump intensity parameter. This compensates
underwriters for unhedgeable gap risk between callbacks.
\end{definition}

\begin{theorem}[Jump Safety]\label{thm:jump-safety}
When $\DeltaPlus$ jumps discretely from $\delta$ to $\delta'$ at a liquidity event
(due to abrupt changes in $\PosCount$ or $\ThetaSum$), the protection value change
is bounded:
\begin{equation}\label{eq:jump-bound}
  |\Delta\text{value}| \leq |V(\delta') - V(\delta)|
\end{equation}
where $V$ is the LP payoff function (\cref{eq:lp-payoff}).
\end{theorem}

\begin{proof}
The LP payoff $V(p)$ is Lipschitz continuous on $[\pl, \pu]$ with
$|V'(p)| = 2p/\pl^2 \leq 2\pu/\pl^2$. The value change at a jump
equals $V(p(\delta')) - V(p(\delta))$, which is bounded by the payoff
difference because $V$ is monotone and no intermediate states are
observable between callbacks. The jump-adjusted funding compensates
for this discontinuity following
\cite[Section~3]{angeris2022perpetuals}. \qedhere
\end{proof}

\subsection{Funding Accumulator}

\begin{definition}[Funding Growth]\label{def:funding-accumulator}
The global funding growth accumulator:
\begin{equation}\label{eq:funding-accumulator}
  g^{\text{funding}} \mathrel{+}= \rfunding_n^{\text{adj}} \cdot \Delta t_n
\end{equation}
This accumulator can be positive or negative. Per-position funding owed
is computed via the Synthetix difference pattern:
$\text{funding}_i = L_i \cdot (g^{\text{funding}} - g^{\text{funding}}_{\text{baseline},i})$.
\end{definition}

\subsection{Convergence Property}

\begin{proposition}[Mark-Index Convergence]\label{prop:convergence}
Under the funding mechanism, $|\pmark - \pindex|$ decreases
after each callback (in expectation), provided $\alpha > 0$.
\end{proposition}

\begin{proof}[Proof sketch]
When $\pmark > \pindex$, PLPs pay positive funding, reducing demand
for the perpetual and pushing $\pmark$ down. When $\pmark < \pindex$,
underwriters pay, increasing demand and pushing $\pmark$ up.
The sensitivity $\alpha$ controls convergence speed.
A formal convergence proof under specific assumptions on arrival rates
is deferred to the Kontrol verification phase. \qedhere
\end{proof}
